%% ch-meta-claim.tex -- Part 2 (\ref{part:meta}), chapter 1: the claim and the trusted base.
%% Cites the design specification as \DS{n} (app-design-spec.tex).
%% Uses only the manual's own preamble; the two macros below are provided here in case this
%% chapter is built before ch-proofs.tex defines its figure macros.
\providecommand{\DS}[1]{\hyperref[ds:#1]{DS-#1}}
\providecommand{\VNBfig}[3]{\[\frac{\displaystyle #2}{\displaystyle #3}\;\;\mbox{\texttt{#1}}\]}
\providecommand{\VNBnohyp}{\phantom{\Gamma}}

\chapter{Self-verification: the claim and the trusted base}
\label{ch:meta-claim}
\index{META}\index{self-verification}

This part describes what VNB does to check its own kernel, what it
proves about it, and what it leaves to hand-written argument or to
assumption.  The work has two purposes (\DS{1}): to make VNB credible as
a proof checker, and to show that code can be modelled formally and its
correctness proven relative to a stated base.

The operations and their checkers are described in
Sections~\ref{sec:kernel-boundary} and~\ref{sec:inference-checked}; the
closed list of files that may record an inference in
Section~\ref{sec:kernel-checked}; the proof certificate in
Section~\ref{sec:proof-certificate}.  They are not repeated here.

%% =========================================================
\section{The claim}\label{sec:meta-the-claim}

\subsection{Consistency is relative}\label{sec:meta-relative}

VNB cannot prove its own consistency, and no check in this part can
either.  This is G\"odel's second incompleteness theorem~\cite{Godel1931}:
a consistent theory with effectively listed axioms, strong enough to
express arithmetic, does not prove its own consistency.  VNB is such a
theory.

So no statement of the form ``VNB cannot prove \textsc{falsity}'' is made
(\DS{2}).  Every consistency statement is relative.  A derivation of
\textsc{falsity}, once found, is conclusive.  Finding none is evidence,
never proof.

\subsection{Version-X}

Every claim is relative to a fixed class theory, called
\emph{version-X}.  Version-X is NBG with global choice (\DS{3}).  Global
choice is needed because \texttt{choice-axiom} gives
$\mathrm{CHOICE}(A) \in A$ for every inhabited class $A$.
$\mathrm{CHOICE}(A)$ denotes exactly when $A$ is inhabited (\DS{52}).  The
definedness certificate never accepts a $\mathrm{CHOICE}$ term by its
shape: it is not a total head.

What VNB calls a class is read as a class of version-X (\DS{4}).  The
oracles are excluded from the formal reasoning and assumed correct
(\DS{5}).

NBG proves the same statements about sets as ZF, and is consistent
exactly when ZF is.  The claim therefore says that VNB is as safe as
ordinary set theory, under the assumptions listed below.

\subsection{The statement}

The claim, in the form this part establishes it, is:
\begin{center}
\emph{Relative to version-X, and assuming the oracles correct, every
inference the kernel accepts is admissible over the core.}
\end{center}
``Admissible over the core'' is defined in
Chapter~\ref{ch:meta-k2}: whenever the hypotheses of the inference are
derivable by the rules of a small, fixed system (the \emph{core}), so is
its conclusion.

\subsection{From the claim to relative consistency}\label{sec:meta-chain}

The step from the claim to relative consistency has five links.  Only
the first is a VNB theorem.
\begin{enumerate}
\item \emph{Admissibility} (formal).  Each non-oracle kernel rule is
  admissible over the core.  This is K2 (\DS{8}), stated rule by rule in
  Chapter~\ref{ch:meta-k2}; its coverage is in
  Section~\ref{sec:meta-coverage}.
\item \emph{Faithfulness} (hand review).  Each transcribed relation is the
  relation its Scheme checker enforces.  This is K3 (\DS{9}).
\item \emph{Soundness of the core} (hand-written).  Each core rule is
  sound for version-X under the reading of classes of \DS{4} and the
  reading of functoids of Chapter~\ref{ch:meta-functoids}; each axiom of
  the catalog holds under that reading (\DS{6}).  The argument is
  Chapter~\ref{ch:meta-soundness}.
\item \emph{Oracles} (assumption, \DS{5}).
\item \emph{Infrastructure} (assumption).  The expression layer, the write
  point and grounding, the install, the certificate store and MIT Scheme
  behave as described.
\end{enumerate}

\begin{prop}[informal; conditional]\label{prop:meta-relcons}
If version-X is consistent, links 1--5 hold, and every installed theorem
was installed by \texttt{qed} over a grounded deduction graph whose every
inference passed its checker, then the installed theorems do not include
\textup{\textsc{falsity}}.
\end{prop}

The proposition is stated to make the dependency explicit.  It is not a
VNB theorem.  What is not yet established for each link is listed in
Chapter~\ref{ch:meta-limits}.

The META theorems are themselves VNB theorems.  They presuppose the
consistency of the system they are proven in, and many of them use the
oracles \texttt{arith} and \texttt{ineq}.

%% =========================================================
\section{The trusted base}\label{sec:meta-base}

\subsection{The write point and the check}

\texttt{dg-apply-rule!} (\texttt{deduction-graphs.scm}) is the only
procedure that writes an inference into a deduction graph.  The
\emph{kernel} consists of all the procedures that call it (\DS{11}).
Before anything is written, \texttt{dg-apply-rule!} hands the proposed inference to the checker
registered for the rule's name (\DS{12}).  A missing checker or a
refusal is an error, and nothing reaches the graph.

Each rule is a relation between the premises of an inference, its
conclusion and its side data.  A checker receives an inference after it
is built and tests whether these stand in the relation of its rule; it
builds nothing itself.  It is written
independently of the procedure that builds the inference, and calls
neither that procedure nor a copy of it.  So the builders and the
tactics carry no soundness: a wrong builder can only fail to prove.

The four checker files register 64 checkers: 40 in
\texttt{rule-checkers-logic.scm}, 13 in \texttt{-schema.scm}, 4 in
\texttt{-rewrite.scm} and 7 in \texttt{-oracle.scm}.  The 57 that are
not oracles are the subject of Chapter~\ref{ch:meta-k2}.

No rule is accepted on trust: the checker registered for each rule
tests the inference, and none accepts every inference.  Two procedures
could register such an always-accepting checker: \texttt{dg-check-arm-with!}, for a documented
tag without a checker, and \texttt{rco-trust-oracle!}.  The first
registers nothing in a complete load, because every documented tag has
its checker; the second has no caller.  The suite checks both.

\paragraph{One exception: the definedness test is shared.}  Several
rules are sound only when some term denotes.  For instance,
\texttt{forall-elim} passes from $\forall x\,A$ to $A[t/x]$ only for a
defined $t$.  When no hypothesis of the inference says that $t$ is
defined, the checker must decide it, and for this it does not have a
test of its own.  It calls \texttt{chk-elim-certified?}
(\texttt{rule-checkers-logic.scm}), which tries a few clauses and then
falls back on \texttt{pi--defined?}, the same definedness test that the
builders use.

The checkers of \texttt{forall-elim}, \texttt{forsome-intro},
\texttt{if-true}, \texttt{if-false}, \texttt{tuples-intro},
\texttt{union-intro} and \texttt{lambda-type} take this route.  The
checker of \texttt{reflexivity} and the value test of the macete door
call \texttt{pi--defined?} directly.

The clauses of \texttt{pi--defined?} are themselves proven: each is
transcribed as a VNB relation and derived over the core, so a term that
a clause accepts does denote (Section~\ref{sec:meta-lemma-star}).  What
the checkers share with the builders is only the Scheme code that
applies those clauses.  A mistake in that code, one that reports a term
as defined when no clause covers it, would reach the builder and its
checker alike, where elsewhere the checker would catch the builder's
mistake.  Whether Scheme code computes its transcribed relation is the
question K3 asks of every checker (Chapter~\ref{ch:meta-limits}).  A
second definedness
test, written separately for the checkers, would remove the overlap; if
one is adopted, it accepts no more than \texttt{pi--defined?} does
(\DS{32}).

\subsection{How ``every inference is checked'' is enforced}

\begin{itemize}
\item The switch \texttt{*dg-check-inferences?*} is armed at load, after
  the four checker files.  \texttt{VNB\_NO\_RULE\_CHECK=1} is the control
  that turns it off.
\item A refusal prints a line beginning \texttt{;VNB RULE CHECKER
  REFUSED}, also inside \texttt{quietly}, and is recorded.  The load ends
  with \texttt{inference checking: $N$ verified, $M$ refusal(s)}; an
  ordinary load (Section~\ref{sec:meta-gates}) is an error when $M>0$.
\item The suite checks that the switch is on after the load, that more
  than a thousand inferences were verified, that no rule is accepted on
  trust, and that an unregistered tag is refused.
\item Each checker has one accepted and one refused control, and a new
  kernel operation comes with its checker and both controls (\DS{13}).
\end{itemize}

\subsection{The code that may write}

Only the files listed in \texttt{*cert-kernel-caller-files*}
(\texttt{certificates.scm}) may call \texttt{dg-apply-rule!}:
\texttt{primitive-inferences}, \texttt{macetes}, \texttt{library},
\texttt{arith-eval} and the four oracle files.  The fatal
\texttt{kernel-callers-audit} refuses a call, or a use as a value, of
\texttt{dg-apply-rule!} anywhere else (\DS{11};
Section~\ref{sec:kernel-checked}).

\subsection{The axioms and the primitive shelf}

The axioms are assumed.  The set and class formation schemas are kernel
rules, not entries of the axiom table.  The shelf of primitives grows
only by an explicit decision, one fact at a time, with its claim
recorded in the file (\DS{15}).

The stamps \texttt{primitive} and \texttt{definitional} bill nothing.  A
stamp therefore records what it assumes (\DS{15}), and a guarded
definitional equation says nothing off its guards (\DS{31}).

\subsection{The gates that stop a load}\label{sec:meta-gates}

An \emph{ordinary load} is any load of VNB without the variable
\texttt{VNB\_KEEP\_GOING=1}: the prover, the band build and the exam.
That variable selects a development mode, the \emph{keep-going load}, in
which an error in a proof file is recorded and the load goes on, so that
every broken file is listed in one load.  A band built that way has holes
in its library, and the end of the load says so.  (This is unrelated to
\texttt{VNB\_CERTIFIED=strict}, which concerns certificates,
Chapter~\ref{ch:provenance}.)

Among the end-of-load gates, the following stop an ordinary load and are
only reported in a keep-going one:
\begin{itemize}
\item inference checking: any refusal;
\item the \texttt{qed}-failure gate: a \texttt{qed} that failed;
\item the after-\texttt{qed} gate: a command run on a proof that a
  \texttt{qed} had just closed;
\item the guarded-definition gate: a guarded defined head in a
  certificate table or typed in $\mathrm{FUN}(\ldots)$ (\DS{31});
\item the view gate: a view typing installed without its proof (\DS{36});
\item \texttt{connective-arity-audit} and \texttt{constant-binder-audit};
\item the category obligations of \texttt{declare-category!};
\item the predicativity gates for \texttt{comp} and for the binder terms
  (\DS{38});
\item the functoid-hypotheses gate: a violation of H1--H4 or H6
  (\DS{35}; Chapter~\ref{ch:meta-functoids});
\item the library gate;
\item \texttt{kernel-callers-audit}.
\end{itemize}
Under keep-going these gates report and the load continues.  The
free-variable audit, the head-registry sweep and the duplicate audits
warn only.

\subsection{The certified load}\label{sec:meta-certified}

The everyday load installs a theorem from its certificate without
running its proof (Section~\ref{sec:proof-certificate}).  The inferences
of a certified theorem are therefore not re-checked in that load: they
were checked by the run that wrote the record.

The certificate key covers the statement, the definitions of its defined
heads, followed transitively, and every definition or macete a checker
read while accepting an inference (\DS{17}).  The kernel hash covers the
files that may write, \texttt{deduction-graphs}, \texttt{expressions},
\texttt{sequents} and the four checker files; a change to any of them
invalidates every certificate.

The run in which every proof is re-executed is the nightly exam
(Chapter~\ref{ch:checking}).  A kernel change lands only after an exam of
a tree that contains it (\DS{16}).

%% =========================================================
\section{Predicativity: staying within NBG}\label{sec:meta-predicative}

NBG allows class comprehension only over conditions whose quantifiers
range over sets.  The kernel keeps to that (\DS{38}).
\begin{itemize}
\item The membership rules of \texttt{comp}, \texttt{sep},
  \texttt{big-union} and \texttt{vnb-lambda}, and the three induction
  rules, refuse a condition or body with a quantifier not confined to
  sets.  The test reads through defined heads.
\item The definedness certificate treats a \texttt{comp}, \texttt{sep},
  \texttt{big-union} or \texttt{vnb-lambda} term as defined only when its
  condition or body is predicative.
\item \texttt{comp-predicativity-audit} and
  \texttt{binder-predicativity-audit} (\texttt{proof-debt.scm}) list every
  impredicative term named by a proof, an installed formula or a functoid
  body.  Both gates stop an ordinary load.
\item The induction census prints
  \texttt{;; NBG induction census: $n$ impredicative induction(s)}; with
  \texttt{*nbg-induction-fatal?*} on, $n=0$ in any load that completes.
\end{itemize}
The encoded core is NBG-strength in the same way (\DS{39}): each core
schema that forms a class carries the premise $\mathrm{NBG\text{-}OK}$ on
its condition or body (Section~\ref{sec:meta-core}).
